\uuid{QwzT}
\exo7id{4148}
\titre{exo7 4148}
\auteur{quercia}
\organisation{exo7}
\datecreate{2010-03-11}
\isIndication{false}
\isCorrection{false}
\chapitre{Fonction de plusieurs variables}
\sousChapitre{Dérivée partielle}
\cpge{mpsi-mp2i.fonctions-deux-variables.b, pcsi.fonctions-deux-variables.b, ptsi.fonctions-deux-variables.b, mpsi-mp2i.fonctions-deux-variables.c, pcsi.fonctions-deux-variables.c, ptsi.fonctions-deux-variables.c, bcpst1.fonctions-deux-variables.c, bcpst1.fonctions-deux-variables.d, mp-mpi.calcul-differentiel-optimisation.a, pt.fonctions-deux-trois-variables.g}
\module{Analyse}
\niveau{L2}
\difficulte{}

\contenu{
\question{
Trouver ${f_k} : {\R^2} \to \R$ de classe $\mathcal{C}^k$ telle que les $k+1$ dérivées d'ordre $k$ en
$(0,0)$ soient distinctes.
}
}
